% 12/11/2020 - First round of corrections: added Ray's corrections
\documentclass[a4paper, twoside]{report}
\usepackage[utf8]{inputenc}
\usepackage{helvet}
\renewcommand{\familydefault}{\sfdefault}
\usepackage{geometry}
\geometry{
left=16mm,
top=30mm,
right=16mm,
bottom=30mm
}
\usepackage{xcolor}
\definecolor{uga_gray}{rgb}{0.8705,0.8705,0.8705}
\definecolor{uga_dark_gray}{rgb}{0.6,0.6,0.6}

\usepackage[absolute,overlay]{textpos}
\usepackage{graphicx}
\usepackage{array}
\usepackage{multicol}
\usepackage{setspace}
\usepackage{setspace}

\usepackage[subpreambles=true]{standalone}

\newcommand{\PhDname}{Rafael Celestre} 

\newcommand{\thesisTitle}{Investigations of the effect of optical imperfections on partially coherent X-ray beam by combining optical simulations with wavefront sensing experiments}
\newcommand{\thesisTitleFR}{\'Etude de l'effet des imperfections optiques aux faisceaux de rayons X partiellement coh\'erents en combinant les simulations optiques avec les mesures du front d'onde}

\begin{document}

\label{form_first}
%% Numéro National de Thèse (donnée par la bibliothèque à la suite du 1er dépôt)
\newcommand{\NNT}{2020GREXXXXX} 			
%% Nom de l'ED
\newcommand{\ecodoctitle}{\'Ecole Doctorale de Physique} 
%% Numéro de l'école doctorale 
\newcommand{\ecodocnum}{47}
%% Spécialité de doctorat
\newcommand{\PhDspeciality}{Physique de la Matière Condensée et du Rayonnement} 	
%% Établissement de préparation
\newcommand{\PhDworkingplace}{Installation Européenne de Rayonnement Synchrotron (ESRF - the European Synchrotron)}
%% Ville de soutenance
\newcommand{\defenseplace}{Grenoble} 
%% Date de soutenance
\newcommand{\defensedate}{XXX janvier, 2021} 													

%%% Établissement / Institution

\newcommand{\vpostt}{0.1} 																	%% NE PAS MODIFIER / DO NOT MODIFY
\newcommand{\hpostt}{6}																		%% NE PAS MODIFIER / DO NOT MODIFY
\newcommand{\vpos}{0.2}																		%% À modifier au besoin pour aligner le logo verticalement
\newcommand{\hpos}{11.5}																	%% À modifier au besoin pour aligner le logo horizontalement
\newcolumntype{L}[1]{>{\raggedright\let\newline\\\arraybackslash\hspace{0pt}}m{#1}}
\newcolumntype{R}[1]{>{\raggedleft\let\newline\\\arraybackslash\hspace{0pt}}lm{#1}}

%%% JURY

%%% Membre n°1 (Président) / Member n°1 (President)
% \newcommand{\jurynameA}{\textbf{Dr. Vincent \textsc{Favre-Nicolin}}, maître de conférences HDR}
% \newcommand{\juryadressA}{ESRF - the European Synchrotron,\\
% en congé de l'Université Grenoble Alpes}
% \newcommand{\juryroleA}{Président}

\newcommand{\jurynameA}{\textbf{\textbf{Dr. Nome \textsc{Sobrenome}}}, status}
\newcommand{\juryadressA}{Instituto}
\newcommand{\juryroleA}{Président}

%%% Membre n°2 (Rapporteur) / Member n°2 (Reviewer)
\newcommand{\jurynameB}{\textbf{Dr. Chris \textsc{Jacobsen}}, professeur}
\newcommand{\juryadressB}{Dept. of Physics \& Astronomy, Northwestern University,\\Advanced Photon Source, Argonne National Lab., États-Unis}
\newcommand{\juryroleB}{Rapporteur}

%%% Membre n°3 (Rapporteur) / Member n°3 (Reviewer)
\newcommand{\jurynameC}{\textbf{Dr. David \textsc{Paganin}}, professeur adjoint}
\newcommand{\juryadressC}{School of Physics and Astronomy, Monash University, Australie}
\newcommand{\juryroleC}{Rapporteur}

%%% Membre n°4 (Examinateur) / Member n°4 (Examiner)
% \newcommand{\jurynameD}{\textbf{Dr. Christian \textsc{Schroer}}, professeur}
% \newcommand{\juryadressD}{Inst. für Nanostruktur- und Festkörperphysik, Universität Hamburg,\\
% Deutsches Elektronen-Synchrotron DESY, Allemagne}
% \newcommand{\juryroleD}{Examinateur}

\newcommand{\jurynameD}{\textbf{Dr. Nome \textsc{Sobrenome}}, status}
\newcommand{\juryadressD}{Instituto}
\newcommand{\juryroleD}{Examinateur}

%%% Membre n°5 (Directeur de thèse) / Member n°5 (Thesis supervisor)
\newcommand{\jurynameE}{\textbf{Dr. Manuel \textsc{Sanchéz del Río}}, physicien HDR}
\newcommand{\juryadressE}{ ESRF - the European Synchrotron}
\newcommand{\juryroleE}{Directeur de thèse}

%%% Membre n°6 (Co-directeur de thèse) / Member n°6 (Thesis co-supervisor)
\newcommand{\jurynameF}{\textbf{Dr. Thomas \textsc{Roth}}, physicien}
\newcommand{\juryadressF}{ESRF - the European Synchrotron}
\newcommand{\juryroleF}{Co-encadrant de thèse}

%%% Membre n°7 (Invité) / Member n°7 (Guest)
\newcommand{\jurynameG}{\textbf{Dr. Lucia \textsc{Alianelli}}, chercheuse principal }
\newcommand{\juryadressG}{Diamond Light Source Ltd, Royaume-Uni}
\newcommand{\juryroleG}{Examinatrice}

%%% Membre n°8 (Invité) / Member n°8 (Guest)
\newcommand{\jurynameH}{\textbf{Dr. Raymond \textsc{Barrett}}, chef du groupe d'optique des rayons X}
\newcommand{\juryadressH}{ESRF - the European Synchrotron}
\newcommand{\juryroleH}{Invité}

%% Il est possible d'ajouter des membres supplémentaires selon le même modèle / More jury members can be added according to the same model

\label{layout_first}
%%% Mise en page / Page layout      
%%% NE RIEN MODIFIER EXCEPTÉ LA PARTIE CONCERNANT LE JURY (voir \label{jury}) SI BESOIN 
%%%%%%%%%%%%%%%%%%%%%%%%%%%%%%%%%%%%%%%%%%%%%%%%%%%%%%%%%%%%%%%%%%%%%%%%%%%%%%%%%%%%%%%%%%%%%%%%%%%%%%%%%%%%%%%%%%%%%%%%%%%%%%%%%%%%%%%%%%%%%%%%%%%%%%%%%%%%%%%%%%%%%%%

\thispagestyle{empty}
.
\begin{textblock}{5}(0,0)
	\textblockcolour{uga_gray}
	\begin{center}
	    	\includegraphics[scale=0.30]{cover/uga_logo.pdf}
	\end{center}{}
	\vspace{300mm}
\end{textblock}

\begin{textblock}{1}(0.6,9.5)
	
	\Huge{\rotatebox{90}{\color{uga_dark_gray}{\fontsize{38}{54}\selectfont Thèse de doctorat}}}
\end{textblock}

\begin{textblock}{1}(2, 9.5)
	\Large{\rotatebox{90}{\color{uga_dark_gray}{NNT : \NNT}}}
\end{textblock}

\begin{textblock}{1}(\hpos,\vpos)
	\textblockcolour{white}
	\includegraphics[scale=0.1]{cover/uga_logo_2.pdf}	
\end{textblock}


%% Texte
\begin{singlespace}
\begin{textblock}{10}(5.5,1.8)
	\textblockcolour{white}
	\color{black}

	\begin{flushleft}
	    \large \textbf{\textsc{Thèse}} our obtenir le grade de:\\ \bigskip \vfill  
	    \Large \textbf{\textsc{Docteur de l’Université Grenoble Alpes}} \\
        \large Spécialité: \textbf{Physique de la Matière Condensée et du Rayonnement}\\
        \small Arrêté ministériel: 25 mai 2016\\ \bigskip \vfill  
        \large Présentée par\\ \bigskip \vfill 
        \Large \textbf{\textsc{Rafael Celestre}}\\ \bigskip \vfill 
	\end{flushleft}
	\begin{center}
		\Large{\thesisTitle} \bigskip %% Titre de la thèse 
		\vfill
		\large{- ou -} \bigskip
		\vfill
		\Large{\thesisTitleFR} \bigskip %% Titre de la thèse 
		\vfill
	\end{center}
	\begin{flushleft}
		\small Thèse de doctorat dirigée par: \\ \bigskip
		\small
		\begin{tabular}{@{} L{9.5cm} R{4.5cm}}
    		\jurynameE  \\ \juryadressE & \juryroleE \\
    		\jurynameF  \\ \juryadressF & \juryroleF  \\ \bigskip
	\end{tabular} 
		\small préparée à \PhDworkingplace~dans \ecodoctitle~(n$^{\circ}$~\ecodocnum).\\ \bigskip
		\vfill
		Thèse présentée et soutenue publiquement à \defenseplace, le \defensedate~devant le jury composé de :
	\end{flushleft}
	%% Members of the jury
	\small

	% Mettre à jour si des membres ont été ajoutés ou retirés
	\label{jury} 		
	\begin{flushleft}
	\begin{tabular}{@{} L{9.5cm} R{4.5cm}}
		\jurynameA  \\ \juryadressA & \juryroleA \\[5pt]
		\jurynameB  \\ \juryadressB & \juryroleB \\[5pt]
		\jurynameC  \\ \juryadressC & \juryroleC \\[5pt]
		\jurynameD  \\ \juryadressD & \juryroleD \\[5pt]
		\jurynameG  \\ \juryadressG & \juryroleG \\[5pt]
		\jurynameH  \\ \juryadressH & \juryroleH \\[5pt]
	\end{tabular} 
	\end{flushleft}   
% 	\end{center}
\end{textblock}
\end{singlespace}
\end{document}
